%% kapa: exam=1 round=36 year=2025 period=2 form=A subject=accounting subject_ko=회계학 seq=27 number=67
%% source: Q-Net 감정평가사 기출문제 (cst00309) artlSeq=5231525 file=raw/1008/5231525/2025년 제36회 감정평가사 1차시험 2교시 문제지 원본.pdf
%% conversion: kordoc 4.24.51 -> tools/gen1.py (auto); review=proofread-sonnet; flags=none
\begin{question}{67}
(주)감평의 20×1년 기계장치(원가모형 적용)와 관련된 기초 및 기말잔액 자료는 다음과 같다.
\begin{tcolorbox}[enhanced, colback=white, colframe=black, boxrule=0.4pt, arc=0pt, left=4pt, right=4pt, top=4pt, bottom=4pt]
\begin{center}\small
\begin{tabular}{c r r}
계정 & \multicolumn{1}{c}{기초} & \multicolumn{1}{c}{기말} \\ \cmidrule(lr){1-1}\cmidrule(lr){2-2}\cmidrule(lr){3-3}
기계장치 & ￦130,000 & ￦150,000 \\ \cmidrule(lr){1-1}\cmidrule(lr){2-2}\cmidrule(lr){3-3}
감가상각누계액 & (40,000) & (50,000) \\ \cmidrule(lr){1-1}\cmidrule(lr){2-2}\cmidrule(lr){3-3}
손상차손누계액 & (5,000) & (3,000) \\ \cmidrule(lr){1-1}\cmidrule(lr){2-2}\cmidrule(lr){3-3}
장부금액 & ￦85,000 & ￦97,000
\end{tabular}
\end{center}
\end{tcolorbox}
20×1년 중 장부금액 ￦10,000의 기계장치(취득원가 ￦30,000, 감가상각누계액 ￦15,000, 손상차손누계액 ￦5,000)를 처분하고 ￦50,000의 처분이익을 인식하였다. 기계장치와 관련된 (주)감평의 20×1년 순현금유입액(현금유입액-현금유출액)은? (단, 기계장치의 취득과 처분은 현금거래이다.)
\begin{choices}
  \item (-) ￦60,000
  \item (-) ￦30,000
  \item (-) ￦10,000
  \item (+) ￦10,000
  \item (+) ￦30,000
\end{choices}
\end{question}
\answer{4}
